%% ========================================================================
%% Chapter 5: Directory Structure & File Operations on Linux
%% Target: Ubuntu Server 22.04
%% ========================================================================

\chapter{Directory Structure and File Operations on Linux}
\label{ch:filesystem-ops}

This chapter discusses the Linux directory structure based on the
\textit{Filesystem Hierarchy Standard} (FHS), the types of files in Linux
(regular, directory, symlink, and device files), file system navigation via the
command line, advanced file operations related to permissions and ownership
(\texttt{chmod}, \texttt{chown}, \texttt{chgrp}), hard links and symbolic
links, file searching with \texttt{find} and \texttt{locate}, and archiving
and compression (\texttt{tar}, \texttt{gzip}, \texttt{zip}). The module is
designed for Ubuntu Server 22.04 lab sessions with a hands-on approach:
concept, output verification, easy scenarios, and real-world scenarios
commonly encountered by server administrators.

% ------------------------------------------------------------------------
\section{Filesystem Hierarchy Standard (FHS) in Linux}
\label{sec:fhs}

\subsection*{Concept Overview}
Linux organises files and folders in a hierarchical structure rooted at
\texttt{/}. On a server, knowing \textit{"what goes where"} is a basic skill:
configuration lives in \texttt{/etc}, logs are mostly under \texttt{/var/log},
user home directories are in \texttt{/home}, and hardware is represented as
files under \texttt{/dev}. The relevant Unix/Linux principle is
\textit{everything is a file}: not only regular files, but also devices
(\texttt{/dev}) and system/process information (\texttt{/proc}, \texttt{/sys})
are treated like files.

\begin{tipbox}
A useful mindset for server administration:
\textit{``configuration is a file'', ``logs are files'', ``devices are files''}.
Therefore, most troubleshooting ends up at \texttt{cat}, \texttt{less},
\texttt{grep}, and \texttt{ls -l}.
\end{tipbox}

\subsection*{Key directories (brief + function)}
\begin{itemize}
    \item \texttt{/} : root of the entire filesystem.
    \item \texttt{/bin} : essential programs for all users (e.g. \texttt{ls}, \texttt{cp}, \texttt{mv}).
    \item \texttt{/sbin} : essential administration programs (often requires \texttt{sudo}).
    \item \texttt{/etc} : system configuration (generally plain-text files).
    \item \texttt{/var} : frequently changing data (logs, cache, spool). Common: \texttt{/var/log}.
    \item \texttt{/home} : user home directories.
    \item \texttt{/root} : home directory for the \texttt{root} user.
    \item \texttt{/tmp} : temporary files (often cleaned automatically).
    \item \texttt{/usr} : userland applications and libraries (sub: \texttt{/usr/bin}, \texttt{/usr/sbin}).
    \item \texttt{/opt} : third-party/optional applications (often used for custom app deployments).
    \item \texttt{/srv} : data for server services (e.g. web/app).
    \item \texttt{/dev} : \textit{device nodes} (disks, ttys, etc.).
    \item \texttt{/proc}, \texttt{/sys} : \textit{virtual filesystems} from the kernel (process/system status).
    \item \texttt{/mnt}, \texttt{/media} : temporary/removable mount points.
\end{itemize}

\subsection*{Lab 4.1 --- FHS Tour: Explore the System (Easy)}
\textbf{Objective:} identify key directories and distinguish configuration,
logs, and user data.

\begin{enumerate}
    \item Check your current location and display the root contents:
\begin{lstlisting}[language=bash, caption={Viewing the root filesystem and its contents}]
pwd
ls -lah /
\end{lstlisting}

    \item View the contents of the configuration and log directories (briefly):
\begin{lstlisting}[language=bash, caption={Inspecting /etc and /var/log}]
ls -lah /etc | head
ls -lah /var/log | head
\end{lstlisting}

    \item (Optional) Install \texttt{tree} for a visual hierarchy:
\begin{lstlisting}[language=bash, caption={Installing tree and displaying the folder structure}]
sudo apt update
sudo apt install -y tree
tree -L 2 /
\end{lstlisting}
\end{enumerate}

\begin{warningbox}
Do not run \texttt{tree /} without a depth limit on a real server --- the
output can be enormous. Use \texttt{-L} to limit the depth.
\end{warningbox}

\subsection*{Lab 4.2 --- Audit Configuration and Service Log Locations (Real-World)}
\textbf{Scenario:} an admin is asked to check the SSH configuration and look
for access indicators in the logs.

\textbf{Objective:} practise locating standard configuration files and reading
logs safely.

\begin{enumerate}
    \item Verify the SSH configuration file:
\begin{lstlisting}[language=bash, caption={Checking the SSH configuration}]
ls -lah /etc/ssh/
sudo test -f /etc/ssh/sshd_config \
  && echo "sshd_config found" \
  || echo "sshd_config not found"
\end{lstlisting}

    \item Read the first lines of the configuration (without editing):
\begin{lstlisting}[language=bash, caption={Reading configuration safely}]
sudo head -n 30 /etc/ssh/sshd_config
\end{lstlisting}

    \item View recent activity in the authentication log:
\begin{lstlisting}[language=bash, caption={Viewing the authentication log}]
ls -lah /var/log/ | grep -E "auth|secure" || true
sudo tail -n 50 /var/log/auth.log
\end{lstlisting}

    \item (Optional) Filter important events:
\begin{lstlisting}[language=bash, caption={Filtering important events in auth.log}]
sudo grep \
 -E "Failed password|Accepted password|Invalid user" \
 /var/log/auth.log | tail -n 30
\end{lstlisting}
\end{enumerate}

% ------------------------------------------------------------------------
\section{File Types in Linux (Regular, Directory, Symlink, Device)}
\label{sec:filetypes}

\subsection*{Concept Overview}
In Linux, a ``file'' is not just a document. From the output of \texttt{ls -l},
the first character indicates the type:
\begin{itemize}
    \item \texttt{-} : regular file
    \item \texttt{d} : directory
    \item \texttt{l} : symbolic link
    \item \texttt{b} : block device (disk/partition)
    \item \texttt{c} : character device (tty/serial)
    \item \texttt{p} : FIFO (named pipe)
    \item \texttt{s} : socket
\end{itemize}

\subsection*{Commands used}
\begin{itemize}
    \item \texttt{ls -l} / \texttt{ls -lah}: view file types and permissions.
    \item \texttt{stat <file>}: detailed metadata (inode, size, timestamps, permissions).
    \item \texttt{file <file>}: guess the content type (ASCII text, image, ELF binary, etc.).
\end{itemize}

\subsection*{Lab 4.3 --- Reading File Types from ls -l Output (Easy)}
\textbf{Objective:} distinguish regular files, directories, symlinks, and device files.

\begin{enumerate}
    \item Create the lab workspace:
\begin{lstlisting}[language=bash, caption={Creating the workspace for this chapter}]
mkdir -p ~/lab-os/week05
cd ~/lab-os/week05
pwd
\end{lstlisting}

    \item Create example files and directories:
\begin{lstlisting}[language=bash, caption={Creating a regular file and example directory}]
mkdir -p data
echo "hello server" > data/hello.txt
ls -lah
ls -lah data
\end{lstlisting}

    \item Create a simple symlink:
\begin{lstlisting}[language=bash, caption={Creating a symlink}]
ln -s data/hello.txt hello-link.txt
ls -lah
\end{lstlisting}

    \item Inspect metadata:
\begin{lstlisting}[language=bash, caption={Inspecting file metadata}]
stat data/hello.txt | head -n 20
file data/hello.txt
\end{lstlisting}
\end{enumerate}

\subsection*{Lab 4.4 --- Reading Device Files and Virtual FS (Real-World)}
\textbf{Scenario:} an admin needs to distinguish device nodes (block vs char)
and virtual files from the kernel.

\textbf{Objective:} recognise \texttt{/dev} (device nodes), \texttt{/proc} and
\texttt{/sys} (virtual filesystems).

\begin{enumerate}
    \item Identify block devices and character devices:
\begin{lstlisting}[language=bash, caption={Block vs character device}]
ls -l /dev/sd* 2>/dev/null | head -n 10
ls -l /dev/tty
\end{lstlisting}

    \item Show that \texttt{/proc} contains process information (virtual):
\begin{lstlisting}[language=bash, caption={Reading system information from /proc}]
ls -lah /proc | head
cat /proc/cpuinfo | head -n 20
\end{lstlisting}

    \item Use \texttt{file} to distinguish content types:
\begin{lstlisting}[language=bash, caption={Testing content type with file}]
file /dev/tty
file /proc/cpuinfo
\end{lstlisting}
\end{enumerate}

\begin{warningbox}
Do not write arbitrarily to \texttt{/dev/*} or \texttt{/proc/*}. Many entries
are sensitive and can affect the system.
\end{warningbox}

% ------------------------------------------------------------------------
\section{File System Navigation via the Command Line}
\label{sec:navigasi}

\subsection*{Concept Overview}
On Ubuntu Server, the terminal is the primary administration interface. Navigation
skills --- finding file locations, moving between directories, and reading
structures --- are the foundation for all advanced tasks.

\subsection*{Commands used}
\begin{itemize}
    \item \texttt{pwd}: display the current working directory.
    \item \texttt{ls}: list directory contents (common options: \texttt{-l}, \texttt{-a}, \texttt{-h}).
    \item \texttt{cd}: change directory.
    \item \texttt{pushd}/\texttt{popd}: directory stack navigation (convenient for switching back and forth).
\end{itemize}

\subsection*{Lab 4.5 --- Safe and Efficient Navigation (Easy)}
\textbf{Objective:} practise absolute paths, relative paths, and \texttt{pushd/popd}.

\begin{enumerate}
    \item Move to root, then return to home:
\begin{lstlisting}[language=bash, caption={Absolute path vs returning to home}]
cd /
pwd
cd ~
pwd
\end{lstlisting}

    \item Practise relative paths:
\begin{lstlisting}[language=bash, caption={Relative paths in the workspace}]
cd ~/lab-os/week05
mkdir -p a/b/c
cd a/b
pwd
cd ../..
pwd
\end{lstlisting}

    \item Use \texttt{pushd/popd} to switch between directories:
\begin{lstlisting}[language=bash, caption={Directory stack navigation}]
pushd /etc
pwd
popd
pwd
\end{lstlisting}
\end{enumerate}

\begin{warningbox}
The command \texttt{rm -rf} deletes directories and their contents without
confirmation. During lab sessions, use \texttt{rm -ri} or verify the path
carefully first.
\end{warningbox}

\subsection*{Lab 4.6 --- Locating Configuration Files from Service Output (Real-World)}
\textbf{Scenario:} you are asked to find the service unit location and related
configuration file.

\textbf{Objective:} use navigation together with reading common system paths on
a server.

\begin{enumerate}
    \item View SSH status and location hints:
\begin{lstlisting}[language=bash, caption={Tracing service information via systemctl}]
systemctl status ssh --no-pager | head -n 30
\end{lstlisting}

    \item Display the service unit file:
\begin{lstlisting}[language=bash, caption={Finding the service unit file}]
systemctl cat ssh --no-pager | head -n 80
\end{lstlisting}

    \item Navigate quickly to the configuration directory:
\begin{lstlisting}[language=bash, caption={Navigating to the SSH configuration directory}]
cd /etc/ssh
pwd
ls -lah | head
\end{lstlisting}
\end{enumerate}

% ------------------------------------------------------------------------
\section{Advanced File Operations: chmod, chown, chgrp}
\label{sec:perm-owner}

\subsection*{Concept Overview}
Access control in Linux is built from three components:
\begin{itemize}
    \item \textbf{Permissions} for the \textit{owner}, \textit{group}, and \textit{others}
    \item \textbf{Ownership}: \textit{user owner} and \textit{group owner}
    \item \textbf{The execute bit on directories}: on a directory, \texttt{x} means
          ``allowed to enter/traverse the directory''.
\end{itemize}

\subsection*{Reading permissions: rwx for user/group/others}
Example \texttt{ls -l} output:
\begin{center}
\texttt{-rwxr-x---}
\end{center}
Meaning:
\begin{itemize}
    \item Character 1: file type (\texttt{-}, \texttt{d}, \texttt{l}, etc.)
    \item Next 3: \textbf{owner} (user)
    \item Next 3: \textbf{group}
    \item Last 3: \textbf{others}
\end{itemize}

\subsubsection*{Meaning of r/w/x on files vs directories}
\begin{itemize}
    \item \textbf{Regular file}:
    \begin{itemize}
        \item \texttt{r}: allowed to read the file contents
        \item \texttt{w}: allowed to modify the file contents
        \item \texttt{x}: allowed to execute the file as a program/script
    \end{itemize}
    \item \textbf{Directory}:
    \begin{itemize}
        \item \texttt{r}: allowed to list the file names (e.g. \texttt{ls})
        \item \texttt{w}: allowed to create/delete/rename entries in the directory
        \item \texttt{x}: allowed to enter (\texttt{cd}) and access files inside
    \end{itemize}
\end{itemize}

\begin{warningbox}
A directory without \texttt{x} can still have its names listed (if \texttt{r}
is set), but cannot be entered or accessed. Conversely, a directory with
\texttt{x} but without \texttt{r} can be accessed if you know the exact
filename.
\end{warningbox}

\subsection{chmod: Octal Notation}
\label{subsec:chmod-octal}

Each permission is represented as an octal digit (0--7). Each triad
(\textit{user/group/others}) is calculated by summing the bit values:

\begin{center}
\begin{tabular}{ccc}
    \textbf{r} (read) & \textbf{w} (write) & \textbf{x} (execute) \\
    4                 & 2                  & 1                    \\
\end{tabular}
\end{center}

Sum the active bit values to get the octal digit for each triad. Three octal
digits combined form the complete permission, for example \texttt{755} means
user=7, group=5, others=5.

\begin{table}[htbp]
    \centering
    \caption{Octal notation to permission conversion table}
    \begin{tabular}{cll}
        \hline
        \textbf{Digit} & \textbf{Permission} & \textbf{Description} \\
        \hline
        0 & \texttt{-{}-{}-} & No access at all \\
        1 & \texttt{-{}-x} & Execute only (enter directory) \\
        2 & \texttt{-w-} & Write only \\
        3 & \texttt{-wx} & Write and execute \\
        4 & \texttt{r-{}-} & Read only \\
        5 & \texttt{r-x} & Read and execute \\
        6 & \texttt{rw-} & Read and write \\
        7 & \texttt{rwx} & Read, write, and execute (full access) \\
        \hline
    \end{tabular}
\end{table}

Usage examples:
\begin{itemize}
    \item \texttt{chmod 755 file} $\rightarrow$ owner \texttt{rwx} (7), group \texttt{r-x} (5), others \texttt{r-x} (5)
    \item \texttt{chmod 640 file} $\rightarrow$ owner \texttt{rw-} (6), group \texttt{r--} (4), others \texttt{---} (0)
    \item \texttt{chmod 600 file} $\rightarrow$ owner \texttt{rw-} (6), group \texttt{---} (0), others \texttt{---} (0)
\end{itemize}

\subsection*{chmod: symbolic notation}
\begin{itemize}
    \item Target: \texttt{u} (user), \texttt{g} (group), \texttt{o} (others), \texttt{a} (all)
    \item Operator: \texttt{+} (add), \texttt{-} (remove), \texttt{=} (set exactly)
\end{itemize}

Examples:
\begin{itemize}
    \item \texttt{chmod u+x script.sh} (add execute for owner)
    \item \texttt{chmod go-rw file} (remove read+write for group and others)
\end{itemize}

\subsection*{Special bits: setuid, setgid, sticky}
In addition to the three main triads, there are special bits:
\begin{itemize}
    \item \textbf{setuid (4xxx)}: program runs with the file owner's privileges (usually root). Example: \texttt{/usr/bin/passwd}.
    \item \textbf{setgid (2xxx)}:
    \begin{itemize}
        \item on a file: runs with the file's group
        \item on a directory: new files inherit the directory's group (\textit{important for shared folders})
    \end{itemize}
    \item \textbf{sticky bit (1xxx)} on a directory: only the file owner (or root) may delete files in that directory.
          Example: \texttt{/tmp}.
\end{itemize}

\subsection*{Ownership: chown and chgrp}
\begin{itemize}
    \item \texttt{chown user file} changes the user owner
    \item \texttt{chown user:group file} changes both user and group at once
    \item \texttt{chgrp group file} changes the group
\end{itemize}

\begin{tipbox}
Access problems on a server are often not because a file is ``missing'', but
because ownership or permissions are incorrect. Make it a habit to verify with
\texttt{ls -l} and \texttt{stat}.
\end{tipbox}

\subsection*{Lab 4.7 --- Basic Permissions on Files and Directories (Easy)}
\textbf{Objective:} observe the difference in meaning of \texttt{x} on a file
vs a directory, and practise octal conversion.

\begin{enumerate}
    \item Prepare test objects:
\begin{lstlisting}[language=bash, caption={Setting up permission test objects}]
cd ~/lab-os/week05
mkdir -p permtest/dirA
echo "hello" > permtest/fileA.txt
ls -ld permtest permtest/dirA
ls -l permtest/fileA.txt
\end{lstlisting}

    \item Change the file permission to \texttt{640} and verify:
\begin{lstlisting}[language=bash, caption={chmod on a file using octal}]
chmod 640 permtest/fileA.txt
ls -l permtest/fileA.txt
\end{lstlisting}

    \item Change the directory permission to \texttt{700}, then try to access it:
\begin{lstlisting}[language=bash, caption={chmod on a directory and testing access}]
chmod 700 permtest/dirA
ls -ld permtest/dirA
cd permtest/dirA
pwd
cd -
\end{lstlisting}

    \item Experiment: remove \texttt{x} from the directory and observe the effect:
\begin{lstlisting}[language=bash, caption={Testing the effect of removing x from a directory}]
chmod 600 permtest/dirA
ls -ld permtest/dirA
cd permtest/dirA 2>/dev/null \
 || echo "Failed: no execute (x) on the directory"
\end{lstlisting}
\end{enumerate}

\subsection*{Lab 4.8 --- Shared Project Folder + Sticky Bit (Real-World)}
\textbf{Scenario:} a shared project folder where group members can collaborate,
new files inherit the group, and users cannot accidentally delete each other's
files.

\textbf{Objective:} apply \texttt{chgrp}, \texttt{chmod 2xxx}, and sticky bit
on a collaborative directory.

\begin{enumerate}
    \item Create the group and project folder:
\begin{lstlisting}[language=bash, caption={Creating the group and project folder}]
sudo groupadd project 2>/dev/null \
  || echo "group project already exists"
sudo usermod -aG project $USER
sudo mkdir -p /srv/project-gamma
sudo chgrp project /srv/project-gamma
\end{lstlisting}

    \item Apply setgid so new files inherit the group, and restrict others:
\begin{lstlisting}[language=bash, caption={Setgid on the project directory}]
sudo chmod 2770 /srv/project-gamma
ls -ld /srv/project-gamma
\end{lstlisting}

    \item Apply the sticky bit so users cannot delete each other's files:
\begin{lstlisting}[language=bash, caption={Sticky bit on the collaborative directory}]
sudo chmod +t /srv/project-gamma
ls -ld /srv/project-gamma
\end{lstlisting}

    \item Test file creation (note: group membership may require re-login):
\begin{lstlisting}[language=bash, caption={Testing collaborative access}]
touch /srv/project-gamma/test-$USER.txt 2>/dev/null \
  || echo "If it fails: re-login or run newgrp"
ls -l /srv/project-gamma | head
\end{lstlisting}
\end{enumerate}

\begin{exercisebox}[4.1]
Determine the chmod number for:
\begin{enumerate}
    \item File: owner \texttt{rw-}, group \texttt{r--}, others \texttt{---}
    \item Shared directory: owner+group \texttt{rwx}, others \texttt{---}, new files inherit the group (setgid)
    \item Directory like \texttt{/tmp}: everyone can create files, but cannot delete files belonging to other users (sticky bit)
\end{enumerate}
\end{exercisebox}

% ------------------------------------------------------------------------
\section{Hard Links and Symbolic Links}
\label{sec:links}

\subsection*{Concept Overview}
A link is a way to create an ``alternative name'' for a file.
\begin{itemize}
    \item \textbf{Hard link} (\texttt{ln}): points to the same inode. Deleting the original name does not remove the data as long as another hard link exists.
    \item \textbf{Symbolic link} / symlink (\texttt{ln -s}): a small file that points to a target path. If the target is moved or deleted, the symlink becomes \textit{broken}.
\end{itemize}

\subsection*{Commands used}
\begin{itemize}
    \item \texttt{ln}, \texttt{ln -s}: create links.
    \item \texttt{ls -li}: display inode numbers and link counts.
    \item \texttt{stat}: verify inodes and metadata.
\end{itemize}

\subsection*{Lab 4.9 --- Testing Hard Links vs Symlinks (Easy)}
\textbf{Objective:} observe the difference in inodes and the effect of deletion.

\begin{enumerate}
    \item Prepare a file:
\begin{lstlisting}[language=bash, caption={Creating a source file}]
cd ~/lab-os/week05
echo "version1" > target.txt
\end{lstlisting}

    \item Create a hard link and a symlink:
\begin{lstlisting}[language=bash, caption={Creating a hard link and symlink}]
ln target.txt hardlink.txt
ln -s target.txt symlink.txt
ls -li target.txt hardlink.txt symlink.txt
\end{lstlisting}

    \item Delete the original file and check the effect:
\begin{lstlisting}[language=bash, caption={Deleting the target and checking the links}]
rm target.txt
ls -li hardlink.txt symlink.txt
cat hardlink.txt
cat symlink.txt \
  || echo "symlink is broken because the target is gone"
\end{lstlisting}
\end{enumerate}

\begin{warningbox}
Hard links are generally not used for directories (restricted to prevent
structural loops). For a ``shortcut'' to a directory, use a symlink instead.
\end{warningbox}

\subsection*{Lab 4.10 --- Symlink for Application Version Deployment (Real-World)}
\textbf{Scenario:} a common deployment pattern --- separate versioned
application folders with \texttt{current} as a symlink pointing to the active
version.

\textbf{Objective:} understand why symlinks are very useful for quick rollbacks.

\begin{enumerate}
    \item Create the version structure:
\begin{lstlisting}[language=bash, caption={Creating dummy application versions}]
cd ~/lab-os/week05
mkdir -p app/releases/v1 app/releases/v2 app/shared
echo "APP_VERSION=v1" > app/releases/v1/.env
echo "APP_VERSION=v2" > app/releases/v2/.env
\end{lstlisting}

    \item Create a \texttt{current} symlink pointing to the active version:
\begin{lstlisting}[language=bash, caption={Creating the current symlink}]
ln -s releases/v1 app/current
ls -lah app
cat app/current/.env
\end{lstlisting}

    \item Simulate an upgrade and rollback by simply switching the symlink:
\begin{lstlisting}[language=bash, caption={Switching versions via symlink}]
rm app/current
ln -s releases/v2 app/current
cat app/current/.env

rm app/current
ln -s releases/v1 app/current
cat app/current/.env
\end{lstlisting}
\end{enumerate}

% ------------------------------------------------------------------------
\section{File Searching with find and locate}
\label{sec:search}

\subsection*{Concept Overview}
On a server, you often need to search for:
\begin{itemize}
    \item configuration files (\texttt{*.conf}), logs, or large files consuming disk space
    \item files with problematic permissions
    \item files modified within a specific time range (for debugging)
\end{itemize}

\subsection*{Commands used}
\begin{itemize}
    \item \texttt{find}: real-time search (accurate, flexible).
    \item \texttt{locate}: database-based search (fast, requires database update).
\end{itemize}

\subsection*{Lab 4.11 --- Basic find: Name, Type, Size (Easy)}
\textbf{Objective:} search for files by name, type, and size.

\begin{enumerate}
    \item Find all \texttt{.txt} files in the workspace:
\begin{lstlisting}[language=bash, caption={find by name and type}]
cd ~/lab-os/week05
find . -type f -name "*.txt"
\end{lstlisting}

    \item Find files larger than 1 MB (common troubleshooting example):
\begin{lstlisting}[language=bash, caption={find by file size}]
find ~ -type f -size +1M 2>/dev/null | head
\end{lstlisting}
\end{enumerate}

\subsection*{Lab 4.12 --- Real-World find: System Configuration Scan (Advanced/Real-World)}
\textbf{Scenario:} search for \texttt{*.conf} configuration files across the
entire system, suppressing permission errors for clean output.

\textbf{Objective:} use \texttt{find} on the root filesystem and manage errors
with redirection.

\begin{lstlisting}[language=bash, caption={Scanning *.conf without flooding permission errors}]
sudo find / -type f -name "*.conf" 2>/dev/null \
  | head -n 30
\end{lstlisting}

\begin{tipbox}
The pattern \texttt{2>/dev/null} is used very often when scanning the entire
system, because many directories simply cannot be read by a regular user.
\end{tipbox}

\subsection*{Lab 4.13 --- locate + updatedb (Easy)}
\textbf{Objective:} use the index database for fast searches.

\begin{enumerate}
    \item Ensure \texttt{plocate} is available:
\begin{lstlisting}[language=bash, caption={Installing the locate tool}]
sudo apt update
sudo apt install -y plocate \
  || sudo apt install -y mlocate
\end{lstlisting}

    \item Update the database (requires root), then search:
\begin{lstlisting}[language=bash, caption={Updating the locate database and searching}]
sudo updatedb
locate sshd_config | head
\end{lstlisting}
\end{enumerate}

\begin{warningbox}
\texttt{locate} is not always real-time. If a file was just created and has not
yet been indexed, the result may be empty until \texttt{updatedb} is run.
\end{warningbox}

\subsection*{Lab 4.14 --- Auditing Problematic Permissions with find (Real-World)}
\textbf{Scenario:} an admin wants to find files/directories that are ``too
open'' (world-writable), or setuid/setgid files, for a light security audit.

\textbf{Objective:} locate objects that are risky from a permissions
perspective.

\begin{enumerate}
    \item Find world-writable directories in home:
\begin{lstlisting}[language=bash, caption={Finding world-writable directories}]
find ~ -type d -perm -0002 2>/dev/null | head -n 30
\end{lstlisting}

    \item Find world-writable files:
\begin{lstlisting}[language=bash, caption={Finding world-writable files}]
find ~ -type f -perm -0002 2>/dev/null | head -n 30
\end{lstlisting}

    \item (Optional) Find setuid/setgid files:
\begin{lstlisting}[language=bash, caption={Auditing setuid/setgid on the system (requires sudo)}]
sudo find / -type f -perm -4000 2>/dev/null \
  | head -n 30   # setuid
sudo find / -type f -perm -2000 2>/dev/null \
  | head -n 30   # setgid
\end{lstlisting}
\end{enumerate}

\begin{warningbox}
Do not arbitrarily change permissions on system setuid/setgid files. This lab
is for identification purposes only.
\end{warningbox}

% ------------------------------------------------------------------------
\section{File Compression and Archiving (tar, gzip, zip)}
\label{sec:archive}

\subsection*{Concept Overview}
An \textbf{archive} combines multiple files/folders into a single package.
\textbf{Compression} reduces the size of the data. On a server, the common
combination is:
\begin{itemize}
    \item \texttt{tar} to archive a folder
    \item \texttt{gzip} (or \texttt{bzip2}/\texttt{xz}) to compress the archive
    \item \texttt{zip} for cross-OS exchange (Windows/macOS) or specific needs
\end{itemize}

\subsection*{tar: most commonly used flags}
\begin{itemize}
    \item \texttt{-c} : create an archive
    \item \texttt{-x} : extract an archive
    \item \texttt{-t} : list contents (view without extracting)
    \item \texttt{-f <file>} : archive filename (required when using a file)
    \item \texttt{-v} : verbose (show files being processed)
    \item \texttt{-C <dir>} : extract/work from a specific directory
    \item \texttt{-z} : use gzip (\texttt{.tar.gz / .tgz})
    \item \texttt{-J} : use xz (\texttt{.tar.xz}) --- stronger compression, slower
\end{itemize}

Quick examples:
\begin{itemize}
    \item Create: \texttt{tar -czvf backup.tar.gz folder/}
    \item List contents: \texttt{tar -tzvf backup.tar.gz}
    \item Extract: \texttt{tar -xzvf backup.tar.gz -C /destination}
\end{itemize}

\subsection*{gzip: important flags}
\begin{itemize}
    \item \texttt{-k} : keep the original file (do not delete the source)
    \item \texttt{-d} : decompress (gunzip)
    \item \texttt{-1} to \texttt{-9} : compression level (1 fastest, 9 maximum)
\end{itemize}
Note: \texttt{gzip} works on a \textbf{single file}. For folders, usually
\texttt{tar} first, then \texttt{gzip}.

\subsection*{zip/unzip: important flags}
\begin{itemize}
    \item \texttt{zip -r} : recursive for folders
    \item \texttt{zip -9} : maximum compression (slower)
    \item \texttt{unzip -l} : list zip contents without extracting
    \item \texttt{unzip -d <dir>} : extract to a specific directory
\end{itemize}

\subsection*{Lab 4.15 --- Archive the Workspace: tar + gzip (Easy)}
\textbf{Objective:} create a \texttt{.tar.gz}, inspect its contents, then
extract it.

\begin{enumerate}
    \item Create a compressed archive of the workspace:
\begin{lstlisting}[language=bash, caption={Creating a tar.gz}]
cd ~/lab-os
tar -czvf week05-backup.tar.gz week05
ls -lah week05-backup.tar.gz
\end{lstlisting}

    \item List the archive contents without extracting:
\begin{lstlisting}[language=bash, caption={Listing tar.gz contents}]
tar -tzvf week05-backup.tar.gz | head -n 30
\end{lstlisting}

    \item Extract to a test folder:
\begin{lstlisting}[language=bash, caption={Extracting tar.gz to a specific folder}]
mkdir -p ~/lab-os/restore-test
tar -xzvf week05-backup.tar.gz \
  -C ~/lab-os/restore-test
ls -lah ~/lab-os/restore-test/week05 | head
\end{lstlisting}
\end{enumerate}

\subsection*{Lab 4.16 --- System Configuration Backup (Subset) + Validation (Real-World)}
\textbf{Scenario:} before making configuration changes on a server, an admin
creates a backup of a subset of \texttt{/etc} and stores it in a backup folder.

\textbf{Objective:} create a tidy, verifiable archive without flooding
permission errors.

\begin{enumerate}
    \item Prepare the backup folder:
\begin{lstlisting}[language=bash, caption={Creating the backup folder}]
mkdir -p ~/backup
\end{lstlisting}

    \item Create a backup of important configuration and verify the size:
\begin{lstlisting}[language=bash, caption={Backing up important configuration with tar.gz}]
sudo tar -czvf ~/backup/etc-core-$(date +%F).tar.gz \
  /etc/ssh /etc/hosts /etc/netplan /etc/systemd \
  2>/dev/null
ls -lah ~/backup/etc-core-*.tar.gz
\end{lstlisting}

    \item Validate the backup contents (without extracting):
\begin{lstlisting}[language=bash, caption={Validating backup contents}]
tar -tzvf ~/backup/etc-core-$(date +%F).tar.gz \
  | head -n 40
\end{lstlisting}

    \item Test restore to a safe directory (do not overwrite the live system):
\begin{lstlisting}[language=bash, caption={Test restore to a sandbox folder}]
mkdir -p ~/backup/restore-sandbox
tar -xzvf ~/backup/etc-core-$(date +%F).tar.gz \
  -C ~/backup/restore-sandbox
ls -lah ~/backup/restore-sandbox/etc | head -n 20
\end{lstlisting}
\end{enumerate}

\begin{tipbox}
The ``backup of important subset'' pattern is safer than archiving all of
\texttt{/etc}, and is more realistic for everyday administration.
\end{tipbox}

\subsection*{Lab 4.17 --- gzip for a Single File (Easy)}
\textbf{Objective:} understand gzip as a single-file compressor, and the effect
of \texttt{-k} and the compression level.

\begin{enumerate}
    \item Create a simple large test file:
\begin{lstlisting}[language=bash, caption={Creating a large test file}]
cd ~/lab-os/week05
yes "logline" | head -n 200000 > biglog.txt
ls -lah biglog.txt
\end{lstlisting}

    \item Compress with gzip and keep the original file:
\begin{lstlisting}[language=bash, caption={gzip with keep}]
gzip -k -9 biglog.txt
ls -lah biglog.txt biglog.txt.gz
\end{lstlisting}

    \item Decompress:
\begin{lstlisting}[language=bash, caption={Decompress gzip}]
gzip -d -k biglog.txt.gz
ls -lah biglog.txt biglog.txt.gz
\end{lstlisting}
\end{enumerate}

\subsection*{Lab 4.18 --- zip for Cross-OS Distribution (Real-World)}
\textbf{Scenario:} you need to send a backup/restore folder to a Windows or
macOS user.

\begin{enumerate}
    \item Ensure zip/unzip is installed:
\begin{lstlisting}[language=bash, caption={Installing zip/unzip}]
sudo apt update
sudo apt install -y zip unzip
\end{lstlisting}

    \item Create a recursive zip and inspect its contents:
\begin{lstlisting}[language=bash, caption={Creating a zip and listing without extracting}]
cd ~/backup
zip -r -9 restore-sandbox.zip restore-sandbox
unzip -l restore-sandbox.zip | head -n 30
\end{lstlisting}

    \item Extract to a new folder:
\begin{lstlisting}[language=bash, caption={Extracting the zip to the target folder}]
mkdir -p ~/backup/unzip-test
unzip restore-sandbox.zip -d ~/backup/unzip-test
ls -lah ~/backup/unzip-test | head
\end{lstlisting}
\end{enumerate}

% ------------------------------------------------------------------------
\section{Summary}
\label{sec:summary-filesystem}

\begin{itemize}
    \item The Linux structure is rooted at \texttt{/} and follows FHS conventions: configuration in \texttt{/etc}, logs in \texttt{/var/log}, home directories in \texttt{/home}.
    \item Linux has many file types: regular, directory, symlink, and device files (\texttt{/dev}) and virtual filesystems (\texttt{/proc}, \texttt{/sys}).
    \item Terminal navigation (\texttt{pwd}, \texttt{ls}, \texttt{cd}, \texttt{pushd/popd}) is the foundation of server administration.
    \item Access control is built from permissions (\texttt{chmod}) and ownership (\texttt{chown}/\texttt{chgrp}); shared folders are a real-world scenario frequently encountered.
    \item Hard links share an inode; symlinks point to a target path and are ideal for deploy/rollback patterns.
    \item \texttt{find} is flexible for real-time searches; \texttt{locate} is fast but depends on a database.
    \item Simple backups are commonly done with \texttt{tar.gz}; \texttt{gzip} handles single files; \texttt{zip} is convenient for cross-OS exchange.
\end{itemize}

% ------------------------------------------------------------------------
\section{Exercises and Independent Tasks}
\label{sec:exercises-filesystem}

\subsection*{A. Conceptual Exercises}
\begin{exercisebox}[4.A]
Choose 8 directories from the FHS list (\texttt{/etc, /var, /home, /tmp, /usr,
/opt, /srv, /dev, /proc, /sys}) and explain the function of each in 1--2
sentences.
\end{exercisebox}

\begin{exercisebox}[4.B]
Explain the difference between a hard link and a symlink, then give one
realistic use case for each.
\end{exercisebox}

\subsection*{B. Lab Exercises (Easy--Intermediate)}
\begin{exercisebox}[4.C]
Create the folder \texttt{\~/lab-os/week05-exercises} containing:
\begin{itemize}
    \item 3 text files, 1 subfolder, and 1 symlink pointing to one of the files.
\end{itemize}
Show proof with \texttt{ls -lah} and \texttt{ls -li}.
\end{exercisebox}

\begin{exercisebox}[4.D]
Set the permissions of \texttt{report.txt} so that the owner can read and
write, the group can only read, and others have no access. Write the command
and verify with \texttt{ls -l}.
\end{exercisebox}

\begin{exercisebox}[4.E]
Use \texttt{find} to locate the 20 largest files in \texttt{/var/log}. If
permission errors appear, discard them with \texttt{2>/dev/null}. Display the
first 20 lines of output.
\end{exercisebox}

\subsection*{C. Lab Exercises (Real-World)}
\begin{exercisebox}[4.F --- Admin Case 1]
You are asked to create the folder \texttt{/srv/data-share} for a team (group
\texttt{project}) with the following rules:
\begin{itemize}
\item all group members can create and edit files
\item new files inherit the folder's group (setgid)
\item other users cannot access the folder
\item users cannot delete files belonging to other users in that folder (sticky bit)
\end{itemize}
Write the commands and show proof with \texttt{ls -ld}.
\end{exercisebox}

\begin{exercisebox}[4.G --- Admin Case 2]
You are asked to gather troubleshooting evidence:
\begin{itemize}
\item copy the last 100 lines of \texttt{/var/log/auth.log} to a file in your home directory
\item compress that file to \texttt{.gz} without deleting the original
\item archive the evidence folder into a \texttt{tar.gz}
\end{itemize}
Write all commands and display the final \texttt{ls -lah} output.
\end{exercisebox}

\subsection*{Reference (Manual Pages)}
\begin{itemize}
    \item \texttt{man hier}, \texttt{man ls}, \texttt{man stat}, \texttt{man file}
    \item \texttt{man chmod}, \texttt{man chown}, \texttt{man chgrp}, \texttt{man umask}
    \item \texttt{man ln}
    \item \texttt{man find}, \texttt{man locate}, \texttt{man updatedb}
    \item \texttt{man tar}, \texttt{man gzip}, \texttt{man zip}, \texttt{man unzip}
\end{itemize}
